\documentclass[11pt]{article}

\usepackage[margin=1in]{geometry}
\usepackage{amsmath,amssymb}
\usepackage{booktabs}
\usepackage{tabularx}
\usepackage{array}
\usepackage{enumitem}
\usepackage[T1]{fontenc}
\usepackage{hyperref}

\newcommand{\code}[1]{\texttt{#1}}
\newcolumntype{L}{>{\raggedright\arraybackslash}X}
\emergencystretch=3em

\title{\code{cluster\_diagnostics\_3d.py}: cluster statistics and radial\\
distribution functions for 3D AMReX plotfiles}
\author{FHDeX}
\date{}

\begin{document}
\maketitle

\section{Overview}

\code{cluster\_diagnostics\_3d.py} (in \code{FHDeX/python}) is the 3D analog of
\code{cluster\_diagnostics.py}, for plotfiles in which $x$, $y$ and $z$ are all
spatial. It was written for the 3D runs of \code{src\_deankow/interaction}. For each
plotfile it
\begin{enumerate}[nosep]
  \item finds the clusters in the density field and measures each one's mass,
        centre, size and peak density;
  \item matches them to the clusters in the previous plotfile, giving each
        cluster a stable id and logging mergers, splits, new clusters and
        clusters that dissolve;
  \item accumulates two radial distribution functions: $g(r)$ of the density
        field and $g(r)$ of the cluster centres.
\end{enumerate}
It needs only Python~3 and numpy, and reuses the plotfile reader in
\code{stack\_plotfiles\_time.py}. Definitions, options and outputs match the 2D
tool; the differences are:
\begin{center}
\begin{tabularx}{\textwidth}{@{}l L L@{}}
\toprule
 & 2D & 3D \\
\midrule
connectivity & 4 face neighbours & 6 face neighbours \\
cluster centre & $(x_c,y_c)$ & $(x_c,y_c,z_c)$; \code{clusters.txt} gains \code{zc} \\
cell measure & $\Delta x\,\Delta y$ & $\Delta x\,\Delta y\,\Delta z$ \\
cluster-$g(r)$ shell & $\pi((r+\Delta r)^2-r^2)/A$ & $\tfrac{4\pi}{3}((r+\Delta r)^3-r^3)/V$ \\
field $g(r)$ & 2D FFT autocorrelation & 3D FFT autocorrelation \\
labelling & per-cell flood fill & vectorized label propagation \\
\code{-{}-mmin} default & 0.01 & $10^{-4}$, and a warning when most mass above the threshold is in lighter regions \\
\bottomrule
\end{tabularx}
\end{center}
At $128^3$ ($2\cdot10^6$ cells) a frame takes about 2~s and 300~MB.

\paragraph{Usage.}
\begin{verbatim}
python3 cluster_diagnostics_3d.py plt_fv3d_* --num-part 2e7 -o diag3d
python3 cluster_diagnostics_3d.py plt_fv3d_* --num-part 2e7 -o diag3d \
        --every 10
\end{verbatim}

\section{Assumptions and limitations}
\begin{itemize}
  \item \textbf{3D, periodic in all directions, level 0 only.} 2D plotfiles are
        rejected with a pointer to the 2D tool.
  \item \textbf{Density field only.} Both $g(r)$ are computed from the density
        variable (\code{phi0} by default).
  \item \textbf{Touching clusters count as one.} In 3D the region above the
        threshold extends further from a cluster's centre than in 2D, so merges
        register at larger separations.
  \item \textbf{Total mass includes negative densities,} so the mass fraction in
        clusters can slightly exceed 1.
  \item \textbf{Choosing \code{-{}-mmin}.} It is a fraction of the total mass.
        The default, $10^{-4}$, comes from the $128^3$ interaction run
        (\code{\_3d\_noise/plt\_fv3d\_060000}). There the regions above the
        $5\times$ threshold split into noise (1--5 cells, mass
        $\le3\cdot10^{-5}$), an empty gap from $3\cdot10^{-5}$ to $10^{-4}$,
        small clusters ($10^{-4}$--$10^{-3}$), and about 470 clusters of
        $10^{-3}$--$3\cdot10^{-3}$ holding 91\% of the mass. The gap is the
        same at thresholds of 2, 10 and $20\times$ the mean. For other runs use
        about $0.05/n_{\mathrm{clusters}}$; the script warns if most of the mass
        above the threshold is dropped. The 2D tool keeps 0.01.
\end{itemize}

\section{Definitions}

The domain is the periodic box
$\prod_d[x_{d,\mathrm{lo}}, x_{d,\mathrm{lo}}+L_d)$ with cells of volume
$\Delta V=\Delta x\,\Delta y\,\Delta z$, centres $\mathbf{x}_i$ and densities
$\phi_i$. The minimum-image displacement is
$\operatorname{mi}_d(a) = (a+\tfrac12L_d)\bmod L_d - \tfrac12L_d$.

\subsection{Cluster}
With $\bar\phi=\sum_i\phi_i\,\Delta V/(L_xL_yL_z)$ and threshold factor $c$
(\code{-{}-threshold}), a cluster is a face-connected set of cells (six
neighbours, periodic) with $\phi_i>c\,\bar\phi$. It is kept if its mass
$M=\sum_{i\in C}\phi_i\,\Delta V$ is at least \code{-{}-mmin} times the total.

\paragraph{Labelling.} Every cell above the threshold starts with a unique
label. Each cell then repeatedly takes the largest label among itself and its
six neighbours, restricted to cells above the threshold, until nothing
changes. This takes about as many sweeps as the largest cluster's diameter in
cells. The labels are then renumbered $1..n$.

With weights $w_i=\phi_i\,\Delta V$:
\begin{center}
\begin{tabularx}{\textwidth}{@{}lL@{}}
\toprule
quantity & definition \\
\midrule
mass $M$ & $\sum_{i\in C}w_i$ \\
centre $(x_c,y_c,z_c)$ & circular mean per direction,
  $x_c=\frac{L_x}{2\pi}\operatorname{atan2}\bigl(\sum w_i\sin\frac{2\pi x_i}{L_x},\sum w_i\cos\frac{2\pi x_i}{L_x}\bigr)\bmod L_x$ \\
$R_g$ & $R_g^2=\sum w_i|\mathbf d_i|^2/M$, with $\mathbf d_i$ the minimum-image displacement from the centre \\
peak, ncells & $\max_{i\in C}\phi_i$; number of cells \\
wide & 1 if some $|d_{i,k}|\ge0.45\,L_k$ (minimum image ambiguous) \\
\bottomrule
\end{tabularx}
\end{center}

\subsection{Tracking between frames}
These are the 2D rules with 3D minimum-image distances. Let
$O_{ba}=\sum_{i\in b\cap a}\phi^{\mathrm{old}}_i\,\Delta V$.
\begin{enumerate}[nosep]
  \item \textbf{Matching:} old $b$ goes to $\arg\max_aO_{ba}$ if
        $O_{ba}>0.1M_b$, else to the nearest new centre within
        \code{-{}-search-radius}.
  \item \textbf{Merger:} two or more old clusters assigned to one new cluster
        with $M_a\ge0.7\sum M_b$.
  \item \textbf{Split:} $O_{ba}>0.1M_b$ for two or more new $a$.
  \item \textbf{Formed and dissolved:} the remaining unassigned clusters.
  \item \textbf{Ids:} the most massive parent's id; fresh ids otherwise.
\end{enumerate}

\subsection{\texorpdfstring{Density-field $g(r)$}{Density-field g(r)}}
With $n_i=\phi_iN\Delta V$, $M$ cells and $N_{\mathrm{tot}}=\sum n_i$,
\begin{equation}
  G(\boldsymbol\Delta)=\sum_in_in_{i+\boldsymbol\Delta}-\delta_{\boldsymbol\Delta,0}\sum_in_i,
  \qquad g(\boldsymbol\Delta)=\frac{M\,G(\boldsymbol\Delta)}{N_{\mathrm{tot}}(N_{\mathrm{tot}}-1)},
\end{equation}
computed by a 3D FFT; a Poisson field gives $g=1$ everywhere. Displacements
are wrapped componentwise into $[-M_d/2,M_d/2)$, and $g(r)$ in a bin is the
average over the lattice displacements in it, which is an exact per-shell
normalization.

\subsection{\texorpdfstring{Cluster-centre $g(r)$}{Cluster-centre g(r)}}
\begin{equation}
  g(r_k)=\frac{\sum_{\text{frames}}\#\{\text{pairs with } r\in[r_k,r_k+\Delta r)\}}
  {\sum_{\text{frames}}\tfrac12n_c(n_c-1)\,\tfrac{4\pi}{3}\bigl((r_k+\Delta r)^3-r_k^3\bigr)/V},
\end{equation}
with $V=L_xL_yL_z$. The shell volume is exact for $r\le L/2$, because each
sphere lies inside the minimum-image cube.

\subsection{\texorpdfstring{Structure factor $S(k)$}{Structure factor S(k)}}
As in 2D, with 3D wave vectors. From the same cell counts as the field $g(r)$,
\begin{equation}
  S(\mathbf k) = \frac{|\hat n(\mathbf k)|^2}{N_{\mathrm{tot}}}, \qquad
  \hat n(\mathbf k) = \sum_i n_i\,e^{-i\mathbf k\cdot\mathbf x_i}, \qquad \mathbf k\ne0 ,
\end{equation}
so that a Poisson field gives $S=1$ and
$S(k)=1+(N/V)\int(g-1)e^{i\mathbf k\cdot\mathbf r}d\mathbf r$. Wave vectors come
from \code{fftfreq}: index $i>N_d/2$ is the negative wavenumber $i-N_d$, so
$k_d=2\pi\cdot(\text{index})/L_d$. $S$ is averaged over shells
$|\mathbf k|\in[(j-\tfrac12)\Delta k,(j+\tfrac12)\Delta k)$ with
$\Delta k$ = \code{-{}-dk} (default $2\pi/\max L_d$), each divided by its exact
number of modes. Only complete shells up to the smallest Nyquist wavenumber
$\pi/\Delta x_d$ are written, averaged over the frames in \code{-{}-t-range}.
For many samples use the in-code structure factor of
\code{src\_deankow/interaction} (\code{struct\_fact\_int}), which writes the same
shell average.

\subsection{Orientational order: crystal or amorphous}
$g(r)$ averages over directions, so a cluster crystal and an amorphous packing
with the same spacing can have the same $g(r)$ peaks. Each cluster centre $i$
is joined to its \code{-{}-nn} nearest neighbours (default 12, minimum image),
or to all neighbours closer than \code{-{}-nn-cut}, and for $l=4,6$
\begin{equation}
  q_{lm,i} = \frac1{n_i}\sum_j Y_{lm}(\hat{\mathbf r}_{ij}) ,\qquad
  q_{l,i} = \Bigl(\frac{4\pi}{2l+1}\sum_{m=-l}^{l}|q_{lm,i}|^2\Bigr)^{1/2} .
\end{equation}
The local order is $\overline{q_{l,i}}$. The global $Q_l$ uses
$\overline{q_{lm,i}}$, averaged over all clusters before the norm: it equals
the local value for a single crystal of any orientation, and is
$\approx0.4/\sqrt N$ for an amorphous packing or a fine polycrystal. Reference
values $(q_4, q_6)$, for face-centred cubic (fcc), hexagonal close-packed
(hcp), body-centred cubic (bcc) and simple cubic lattices (defined in
\code{order\_diagnostics.tex}): fcc (12 nn) $(0.191, 0.575)$, hcp (12) $(0.097, 0.485)$,
bcc (8) $(0.509, 0.629)$, simple cubic (6) $(0.764, 0.354)$, random points (12)
$(0.28, 0.28)$.

As in 2D, $\mathrm{PR}=(\sum S)^2/(n\sum S^2)$ on the $S(k)$ shell with the
largest mean $S$ is $\approx(\text{number of spots})/n$ for Bragg spots (8/602
for an fcc lattice of blobs) and $\approx0.5$ for a ring.

\section{Options}
\begin{center}
\begin{tabularx}{\textwidth}{@{}l l L@{}}
\toprule
option & default & meaning \\
\midrule
\code{plotfiles} & (required) & 3D plotfile directories \\
\code{-{}-num-part N} & (required) & \code{num\_part} of the run \\
\code{-{}-var NAME} & \code{phi0} & density variable \\
\code{-{}-threshold C} & 5 & cluster cells have $\phi>C\bar\phi$ \\
\code{-{}-mmin F} & $10^{-4}$ & smallest cluster kept, as a fraction of the total mass; about $0.05/n_{\mathrm{clusters}}$ \\
\code{-{}-every K} & 1 & every $K$-th plotfile \\
\code{-{}-rmax R} & $L/2$ & range of both $g(r)$; values above $L/2$ are clipped with a warning \\
\code{-{}-dr D} & $\Delta x$ & field $g(r)$ bin width \\
\code{-{}-dr-clusters D} & $L/50$ & cluster $g(r)$ bin width \\
\code{-{}-search-radius R} & $L/4$ & largest centre jump allowed when matching without overlap \\
\code{-{}-t-range T0 T1} & all & time window for the $g(r)$ averages \\
\code{-{}-rdf-per-frame} & off & also write the per-frame field $g(r)$ \\
\code{-{}-dk D} & $2\pi/\max L$ & shell width of the structure factor \\
\code{-{}-sf-per-frame} & off & also write the per-frame shell-averaged $S(k)$ \\
\code{-{}-nn K} & 12 & bond order: number of nearest neighbours \\
\code{-{}-nn-cut R} & off & bond order: all neighbours closer than $R$ instead \\
\code{-{}-top-modes K} & 12 & \code{order.txt} \code{top\_share}: largest modes on the $S(k)$ peak shell \\
\code{-{}-sf-full} & off & write the full $S(\mathbf k)$, averaged over the window, as the plotfile \code{sf\_full\_avg} \\
\code{-{}-sf-full-per-frame} & off & also write the full $S(\mathbf k)$ of each frame as \code{sf\_full\_<step>} \\
\code{-o DIR} & \code{cluster\_diag\_3d} & output directory \\
\bottomrule
\end{tabularx}
\end{center}

\section{Output files}
These are the same as the 2D tool (see \code{README\_cluster\_diagnostics.md}),
with a \code{zc} column added:
\begin{itemize}[nosep]
  \item \code{summary.txt}: \code{t step nclusters mass\_fraction largest\_mass merges splits forms dissolves}
  \item \code{clusters.txt}: \code{t step id mass xc yc zc Rg peak ncells wide q4 q6}
  \item \code{events.txt}: \code{t step type old\_ids old\_masses new\_ids new\_masses}
  \item \code{rdf\_field.txt}: \code{r g ndisplacements} (and per-frame columns)
  \item \code{rdf\_clusters.txt}: \code{r\_lo r\_hi g pairs expected}
  \item \code{sf.txt}: \code{k S nmodes} (and per-frame columns): the shell-averaged structure factor
  \item \code{order.txt}: \code{t step nclusters q4 q6 Q4\_global Q6\_global}
        \code{kpeak nmodes PR top\_share}
  \item \code{sf\_full\_avg}, \code{sf\_full\_<step>} (with \code{-{}-sf-full},
        \code{-{}-sf-full-per-frame}): the full $S(\mathbf k)$ as a plotfile on
        the 3D $k$ grid, $k=0$ at the centre, components \code{S} and
        \code{log10S} (33~MB at $128^3$)
\end{itemize}

\section{Checks}
\begin{center}
\begin{tabularx}{\textwidth}{@{}L L@{}}
\toprule
test & result \\
\midrule
Poisson field, $64^3$, $N=4\cdot10^6$ & $g=1.00000$, rms deviation $10^{-5}$ \\
$4\times4\times4$ lattice of Gaussian blobs, sites on the boundaries & 64 clusters at the sites ($x_c=1\equiv0$); cluster $g(r)$ only at $L/4$, $\sqrt2L/4$, $\sqrt3L/4$, $L/2$ with pair counts 192, 384, 256, 96 \\
shifts by half a cell and half the box & identical $g(r)$ \\
\code{-{}-rmax 0.7} & clipped to 0.5 with a warning \\
blobs merging across the $z$ boundary, and the reverse & one merge and one split; $0.594+0.394\to0.988$ \\
2D pattern extruded in $z$, against the 2D tool & same 4 clusters, masses and centres \\
3D interaction run, $128^3$, default \code{-{}-mmin} & about 476 clusters (92\% of the mass), mostly 0.001--0.003; $g(r)$ peak at 0.14--0.16; about 2~s per frame \\
sc, bcc, fcc, hcp point lattices; random points & reference $q_4$, $q_6$ above, global $=$ local, rotation invariant; random: $q_6=0.283$, $Q_6=0.017$ ($N=500$) \\
fcc lattice of blobs, $48^3$; random blobs & PR $=0.0133=8/602$; random PR 0.55 \\
\code{\_3d\_noise}, $t=0.032$--$0.038$ & $q_4=0.13$, $q_6=0.41$--$0.42$, $Q_6=0.06$--$0.07$; PR 0.32 over 1142 modes: amorphous (2D \code{\_new\_repul\_2}: $\psi_6=0.99$ local and global, 6 spots) \\
\bottomrule
\end{tabularx}
\end{center}

\end{document}
